% =====================================================================
\BookChapter[Julia]{Julia: One Language for Prototype and Production}
\label{ch:julia}

\begin{researchquestion}
Can a language be as easy to write as Python and as fast as C?
\end{researchquestion}
\OrientationNote
% =====================================================================

\begin{factsheet}{Julia}
First appeared & 2012 (version 1.0 in 2018; version 1.13 in 2026)\\
Designed by & Jeff Bezanson, Stefan Karpinski, Viral B.\ Shah and Alan Edelman, MIT\\
Paradigms & multiple dispatch, functional, array, procedural; metaprogramming (macros)\\
Typing & dynamic, with optional type annotations and a rich type system\\
Execution & just-in-time compiled to native code (via LLVM), specialised per argument type\\
File extension & \file{.jl}\\
Cost & free and open source (MIT licence)\\
Used for & scientific computing, differential equations, optimisation, scientific machine
  learning, climate and ocean modelling, power systems, pharmacometrics, astronomy\\
Famous software & DifferentialEquations.jl, JuMP (optimisation), Flux.jl, ModelingToolkit.jl,
  Makie.jl, Oceananigans.jl (climate), PowerModels.jl\\
\end{factsheet}

Julia was created to solve what its authors called the \emph{two-language problem}.
Scientists like to prototype in a convenient dynamic language (\py, MATLAB, R), but when
the code must run fast they rewrite the core in C, C++ or Fortran. This is slow, error-prone
and creates a barrier between the people who use code and the people who write it. Julia
aims to be one language that is as easy to write as \py or MATLAB and as fast as C or
Fortran. It largely succeeds.

\LearningOutcomes{Explain the two-language problem; read Julia syntax, including Unicode names and the broadcast dot; explain why loops are fast in Julia and why the first call of a function is slow; define methods for your own types with multiple dispatch; and name the main Julia packages for engineering.}

\section{History and philosophy}

\begin{history}
In 2009 four programmers --- Jeff Bezanson, Stefan Karpinski, Viral Shah and Alan
Edelman, a professor of applied mathematics at MIT --- started working on a new language.
Their 2012 blog post ``Why We Created Julia'' began: \emph{``We are greedy: we want more.''}
They wanted the speed of C, the dynamism of Ruby, the mathematical notation of MATLAB, the
general-purpose usefulness of \py, the statistics of R and the macros of Lisp. Julia 1.0,
with a promise of a stable language, was released in 2018. In 2019 the creators received
the Wilkinson Prize for Numerical Software.
\end{history}

Two ideas make Julia work:
\begin{enumerate}
\item \textbf{Just-in-time specialisation.} The first time a function is called with
  arguments of particular types (say, two \code{Float64}s), Julia compiles a version of the
  function specialised for exactly those types, using the LLVM compiler infrastructure
  that also powers Clang and Rust. Later calls run that native code directly. Loops are
  therefore fast; you do not need to vectorise to get speed.
\item \textbf{Multiple dispatch.} A function can have many \emph{methods}, and Julia
  chooses which one to run based on the types of \emph{all} its arguments. This turns out
  to be a very natural way to organise mathematical code (Section~\ref{sec:dispatch}).
\end{enumerate}

\section{Getting started}

Install Julia with \textbf{juliaup} from \url{https://julialang.org}. Start the
interactive REPL by typing \code{julia}, or use VS Code with the Julia extension, Jupyter
(the ``Ju'' in Jupyter stands for Julia) or Pluto.jl, a reactive notebook in which
dependent cells update automatically.
\begin{lstlisting}[language=Julia]
julia> println("Hello, engineers!")
Hello, engineers!

julia> 2 + 3
5

julia> sqrt(80 * 9.81 / 0.25)
56.02856414093976
\end{lstlisting}
The REPL has helpful modes: type \code{?} for help on a function, \code{]} for the package
manager (e.g.\ \code{] add Plots}), and \code{;} for a shell.

\section{The language in brief}

Julia's syntax will look familiar after MATLAB and \py: blocks end with \code{end},
indices start at 1, and comments start with \code{\#}.
\begin{lstlisting}[language=Julia]
g = 9.81                     # Float64
n = 40                       # Int64
material = "S275"            # String
ok = true                    # Bool
z = 3 + 4im                  # Complex
r = 1//3 + 1//4              # Rational: 7//12 (exact)
x = [1.0, 2.0, 3.0]          # Vector{Float64} (a column vector)
A = [4 -2 1; -2 4 -2; 1 -2 4]  # Matrix{Int64}, MATLAB-style literal
t = 0:0.5:20                 # a range (end included), stored compactly
x[1], x[end], x[2:3]         # indexing from 1, `end` as in MATLAB
\end{lstlisting}
Julia accepts Unicode names, so code can look like the equations in a textbook. In the REPL
or VS Code, type \code{\textbackslash alpha} and press Tab to get $\alpha$:
\begin{lstlisting}[language=Julia]
σ_y = 275e6                  # yield stress, Pa (type \sigma<Tab>_y)
ε = σ_y / 210e9              # yield strain
2π                           # juxtaposition: 2*pi = 6.283185307179586
\end{lstlisting}
(Writing a number directly in front of a variable, as in \code{2$\pi$} or \code{3x}, means
multiplication.)

\subsection*{Functions can be written the way mathematics writes them}

\begin{lstlisting}[language=Julia]
# Full form
function reynolds(ρ, v, D, μ=1.0e-3)
    return ρ * v * D / μ
end

# Short "mathematical" form for one-line functions
area(d) = π * d^2 / 4

# Anonymous function
f = x -> x^3 - 2x - 5
\end{lstlisting}

\subsection*{A dot applies any function element by element}

Julia makes element-wise operations explicit with a dot. Any function, including your own,
can be applied element by element by writing a dot after its name, and operators get a dot
in front:
\begin{lstlisting}[language=Julia]
d = [0.05, 0.08, 0.10]
area.(d)                     # apply area to each element
sin.(t) .+ 2 .* t.^2         # element-wise expression
@. sin(t) + 2t^2             # the @. macro adds all the dots for you
\end{lstlisting}
Unlike MATLAB, where only built-in operators have element-wise versions, the dot syntax
works for every function. And unlike NumPy, a chain of dotted operations is \emph{fused}
into a single loop without temporary arrays, which makes it very fast.

\section{Case study: the falling skydiver}

\codefile{Julia}{julia/drag.jl}{Euler solution of the drag equation}
\begin{lstlisting}[style=shell]
$ julia drag.jl
\end{lstlisting}
\outfile[title={\footnotesize\sffamily\color{muted}Output (Julia 1.13)}]{julia/drag.out}
The output is identical, character for character, to the table of the \py, C and
Fortran chapters: the program performs the same double-precision operations in the same
order. The function \code{euler} is generic: it takes
the right-hand side \code{f} as an argument, like SciPy's solvers. \code{@printf} is a
\term{macro} (macros start with \code{@}); it generates fast formatting code at compile
time. The \code{const} declarations tell the compiler that these global values will not
change, which matters for performance.

\section{Speed: loops are fine}

In \py or MATLAB you are told to avoid loops. In Julia a plain loop is as fast as C:
\begin{lstlisting}[language=Julia]
function sum_of_squares(x)
    s = zero(eltype(x))      # zero of the right type (0.0 for Float64)
    for xi in x
        s += xi^2
    end
    return s
end

x = rand(10^7)
@time sum_of_squares(x)      # first call: includes compilation
@time sum_of_squares(x)      # second call: compiled machine code, a few milliseconds
\end{lstlisting}
The first call takes longer because Julia compiles the function. This compilation delay,
sometimes called ``time to first plot'', has been greatly reduced in recent versions by
caching compiled code, but it is still noticeable when loading large packages. Julia is
therefore best for work sessions in which you run many calculations, rather than for tiny
scripts started thousands of times.

\section{Types and multiple dispatch}
\label{sec:dispatch}

You can define your own types with \code{struct} and give a function different methods
for different argument types:
\begin{lstlisting}[language=Julia]
abstract type Section end

struct Rectangle <: Section     # <: means "is a subtype of"
    b::Float64
    h::Float64
end

struct Circle <: Section
    d::Float64
end

area(s::Rectangle) = s.b * s.h
area(s::Circle)    = π * s.d^2 / 4
second_moment(s::Rectangle) = s.b * s.h^3 / 12
second_moment(s::Circle)    = π * s.d^4 / 64

sections = [Rectangle(0.05, 0.10), Circle(0.08)]
for s in sections
    println(typeof(s), ": A = ", area(s), ", I = ", second_moment(s))
end
\end{lstlisting}
\outfile[title={\footnotesize\sffamily\color{muted}Output (Julia 1.13)}]{julia/sections.out}
The numbers agree with the C++ program of Chapter~\ref{ch:cpp}. \code{println} shows
every significant digit, so the rounding error of Section~\ref{sec:floatingpoint} becomes
visible: $0.05 \times 0.10$ is stored as \code{0.005000000000000001}, because neither factor
is exact in binary. \py prints exactly the same digits.
This looks like the C++ example of Chapter~\ref{ch:cpp}, but the methods are not owned by
the classes: \code{area} is one function with several methods, and anyone can add a new
method for a new type, or a new function for existing types, without modifying the
original code. Dispatch on \emph{several} arguments is what makes this powerful: for
example, Julia's \code{*} has different methods for (scalar, matrix), (sparse matrix,
vector), (diagonal matrix, diagonal matrix) and so on, each with the best algorithm, and
the right one is chosen automatically. Packages written by different people combine
remarkably well because of this.

\section{Packages for engineers}

Julia's package manager is built in (\code{using Pkg; Pkg.add("Plots")}), and every project
can record its exact package versions in \file{Project.toml} and \file{Manifest.toml},
which makes results reproducible. Highlights for engineering:
\begin{itemize}
\item \textbf{DifferentialEquations.jl}: one of the most complete ODE/SDE/DAE solver suites
  in any language;
\item \textbf{ModelingToolkit.jl}: equation-based, acausal modelling of physical systems
  (similar to Modelica), with automatic simplification and code generation;
\item \textbf{JuMP}: a modelling language for mathematical optimisation (linear,
  mixed-integer, nonlinear), widely used in power systems and operations research;
\item \textbf{Plots.jl} and \textbf{Makie.jl}: plotting; \textbf{DataFrames.jl}: tables;
\item \textbf{Unitful.jl}: physical units checked at compile time with no run-time cost;
\item \textbf{ForwardDiff.jl}, \textbf{Enzyme.jl}: automatic differentiation of Julia
  code;
\item \textbf{CUDA.jl}: running Julia code on NVIDIA GPUs.
\end{itemize}
Here is the drag problem solved with DifferentialEquations.jl:
\begin{lstlisting}[language=Julia]
using DifferentialEquations, Plots
m, c, g = 80.0, 0.25, 9.81
f(v, p, t) = g - c / m * v^2
prob = ODEProblem(f, 0.0, (0.0, 20.0))
sol = solve(prob, Tsit5(), reltol=1e-10, abstol=1e-10)   # both tolerances
sol(10.0)                      # 52.7496... (the solution is a continuous function)
plot(sol, xlabel="t (s)", ylabel="v (m/s)", label="v(t)")
\end{lstlisting}

\section{Strengths and weaknesses}

\begin{center}
\small
\begin{tabularx}{\linewidth}{@{}>{\raggedright\arraybackslash}X>{\raggedright\arraybackslash}X@{}}
\toprule
\textbf{Strengths} & \textbf{Weaknesses}\\
\midrule
Speed of compiled languages with dynamic-language convenience & Compilation delay on first use\\
Mathematical syntax (Unicode, MATLAB-like arrays) & Smaller community and job market than \py\\
Multiple dispatch: highly composable packages & Some packages less mature or less documented\\
Excellent differential-equation and optimisation libraries & Error messages involving types can be long\\
Free; built-in reproducible package management & Fewer engineers know it, so sharing code can be harder\\
\bottomrule
\end{tabularx}
\end{center}

\section{Further reading}
\begin{itemize}
\item The official documentation and learning resources, \url{https://julialang.org/learning/}.
\item B.~Lauwens and A.~B.~Downey, \emph{Think Julia}, O'Reilly, 2019 (free online).
\item The ``Modern Julia Workflows'' guide, \url{https://modernjuliaworkflows.github.io}.
\item C.~Rackauckas, \emph{Parallel Computing and Scientific Machine Learning} (MIT
  course notes, free online).
\end{itemize}

\begin{exercises}
\item Install Julia and run \file{drag.jl}. Compare the output with the \py table.
\item Write a function \code{newton(f, df, x0; tol=1e-12)} that implements Newton's method
  and test it on $x^3 - 2x - 5 = 0$. (The semicolon introduces \emph{keyword} arguments.)
\item Time \code{sum\_of\_squares} and \code{sum(x.\^{}2)} for $10^7$ random numbers. Then time
  the equivalent pure-\py loop and the NumPy expression. Comment on the results.
\item Add a \code{Tube} type to the cross-section example and define \code{area} and
  \code{second\_moment} methods for it. Did you have to change any existing code?
\item Use \code{Unitful.jl} to compute the terminal velocity with units attached, and
  check that the result is in m/s. What happens if you accidentally give the mass in grams?
\end{exercises}
